% ============================================================
% SCHEDE PG EXTRA — CASA SPINOLA
% @rpg-schema: <https://rpg-schema.org/instances/genova1507/crew-spinola>
%   fitd:hasOptionalMember
%     <.../pg-fra-ambrogio> ,
%     <.../pg-ser-drago> .
% ============================================================

\chapter{Schede PG Extra --- Casa Spinola}
\index{schede PG!Spinola!extra}

% ════════════════════════════════════════════════════════════
% PG EXTRA 1 — FRÀ AMBROGIO
% @rpg-schema: <https://rpg-schema.org/instances/genova1507/pg-fra-ambrogio>
%   a fitd:Character ;
%   rdfs:label "Frà Ambrogio da Savona" ;
%   fitd:crew <.../crew-spinola> ;
%   fitd:role "Il Cappellano / Il Nodo tra Fede e Finanza" ;
%   fitd:action [ fitd:actionName "Study" ; fitd:dots 3 ] ;
%   fitd:action [ fitd:actionName "Consort" ; fitd:dots 4 ] ;
%   fitd:action [ fitd:actionName "Sway" ; fitd:dots 2 ] .
% ════════════════════════════════════════════════════════════

\pgheader[SpinolaRed]{Frà Ambrogio da Savona}{Il Cappellano / Il Nodo tra Fede e Finanza}{Spinola}
\pgquote{La Chiesa perdona i debiti spirituali. La famiglia perdona i debiti finanziari. Io gestisco entrambi.}

\section*{Background \& Motivazione}

\begin{rulebox}{}
  \textbf{Background:} Cinquantasei anni, agostiniano, cappellano di casa
  Spinola da ventitré anni. Non è un uomo di Dio nel senso che Fra' Matteo
  dei Doria è un uomo di Dio — è un amministratore ecclesiastico che ha
  capito prima di molti che le banche e le chiese servono la stessa funzione
  sociale: gestire l'incertezza del futuro. Conosce i bilanci del banco meglio
  di chiunque tranne Raffaele, e conosce i segreti confessionali di metà dei
  debitori della famiglia.

  \medskip\noindent\textcolor{SpinolaRed}{\textbf{Motivazione:}} Frà Ambrogio
  crede — sinceramente, senza cinismo — che una Genova libera sia più adatta
  alla prosperità ordinata del banco. I francesi portano instabilità, e
  l'instabilità è nemica del credito a lungo termine. Questa è la sua teologia.
\end{rulebox}

\section*{Attributi \& Azioni}

\begin{multicols}{2}

\begin{attributoblock}[SpinolaRed]{INSIGHT --- Mente}
  \azione{Caccia}{Hunt}{1}{Trova informazioni attraverso reti parrocchiali}
  \azione{Studio}{Study}{3}{Teologia morale, diritto canonico, contabilità ecclesiastica}
  \azione{Osserva}{Survey}{2}{Legge le persone attraverso ciò che tacciono}
  \azione{Ingegno}{Tinker}{2}{Sistemi di comunicazione sicura, documenti in codice}
\end{attributoblock}

\columnbreak

\begin{attributoblock}[SpinolaRed]{PROWESS --- Corpo}
  \azione{Finezza}{Finesse}{1}{Calligrafia impeccabile, documenti di ogni tipo}
  \azione{Ombra}{Prowl}{1}{Si muove tra conventi e chiese senza essere notato}
  \azione{Combattimento}{Skirmish}{0}{Non combatte}
  \azione{Forza bruta}{Wreck}{0}{No}
\end{attributoblock}

\end{multicols}

\begin{attributoblock}[SpinolaRed]{RESOLVE --- Spirito}
  \begin{multicols}{2}
    \azione{Percezione}{Attune}{2}{Sente la colpa e la paura altrui come un barometro}
    \azione{Comando}{Command}{2}{Autorità quieta, mai imposta — eppure efficace}
    \azione{Relazioni}{Consort}{4}{La rete ecclesiastica genovese è sua quanto dei Fieschi}
    \azione{Persuasione}{Sway}{2}{Convince con la logica della coscienza}
  \end{multicols}
\end{attributoblock}

\medskip
\begin{pgclockbox}{Stress \hfill \clock{9}{0}}
  \small\textit{Traumi: $\square$ Freddo \quad $\square$ Duro \quad $\square$ Ossessionato \quad $\square$ Paranoico \quad $\square$ Irrazionale \quad $\square$ Spezzato}
\end{pgclockbox}

\section*{Abilità Speciale}

% @rpg-schema: <.../pg-fra-ambrogio>
%   fitd:specialAbility [ rdfs:label "Il Peso della Coscienza" ] .

\begin{spinolabox}{Il Peso della Coscienza}
  Frà Ambrogio può dichiarare che un PNG — debitore, funzionario, mercante —
  ha una questione irrisolta di coscienza che lo rende vulnerabile alla
  sua influenza. Il PNG esegue una richiesta moderata per evitare che la
  cosa venga resa pubblica, anche senza minaccia esplicita.

  \medskip\noindent\textit{Una volta per sessione. Richiede che il PNG
  abbia plausibilmente qualcosa da nascondere — il GM conferma.
  Non funziona su personaggi completamente privi di scrupoli.}
\end{spinolabox}

\section*{Legami}

\begin{table}[h]
\centering\renewcommand{\arraystretch}{1.4}
\begin{tabularx}{\linewidth}{@{}L{2.6cm} X@{}}
  \toprule\rowcolor{SpinolaRed}
  \textcolor{white}{\textbf{PG}} & \textcolor{white}{\textbf{Legame}} \\
  \midrule
  Raffaele & Sa che ho letto i libri veri del banco. Siamo complici da
            vent'anni. Questo è più solido della lealtà. \\
  \rowcolor{SpinolaRed!8}
  Lucrezia & La osservo con preoccupazione genuina. Porta troppo.
            Ho provato a dirglielo. Non ha ascoltato. \\
  Ser Benedetto & Siamo la stessa cosa in abiti diversi. Lui lo sa.
                 Nessuno dei due lo dice ad alta voce. \\
  \rowcolor{SpinolaRed!8}
  Ippolito & Lo tratto come il nipote che non ho avuto. Questo mi rende
            cieco ai suoi difetti — e lui lo sa e ne approfitta. \\
  \bottomrule
\end{tabularx}
\end{table}

\section*{Equipaggiamento}
\begin{goldlist}
  \item Breviario con inserti di carta cifrata — corrispondenza tra conventi
  \item Copia non ufficiale del registro dei debiti di Raffaele (solo i nomi ecclesiastici)
  \item Lettera di presentazione firmata dall'Abate di Sant'Andrea — apre ogni porta religiosa
  \item Una piccola reliquia --- autentica. Vale più di qualsiasi documento.
\end{goldlist}

\clearpage

\begin{secretbox}
  \textbf{Segreto:} Frà Ambrogio ha scoperto — attraverso una confessione
  indiretta, dunque senza il vincolo del sigillo — che il Monsignore Gregorio
  Fieschi ha trattato con i francesi. Lo sa da un mese. Non lo ha detto agli
  Spinola perché sta valutando se questa informazione valga di più usata come
  leva su Gregorio o consegnata alla famiglia.

  \medskip\noindent\textcolor{SecretRed}{\textbf{Agenda:}}
  \textit{Usare questa informazione al momento giusto per massimizzare
  l'influenza degli Spinola sulla rete ecclesiastica genovese dopo la rivolta.
  Non per denunciare Gregorio — per farlo lavorare per la famiglia.}
\end{secretbox}

\clearpage

% ════════════════════════════════════════════════════════════
% PG EXTRA 2 — SER DRAGO
% @rpg-schema: <https://rpg-schema.org/instances/genova1507/pg-ser-drago>
%   a fitd:Character ;
%   rdfs:label "Ser Drago detto il Silenzioso" ;
%   fitd:crew <.../crew-spinola> ;
%   fitd:role "Il Braccio / Il Problema Risolto" ;
%   fitd:action [ fitd:actionName "Hunt" ; fitd:dots 3 ] ;
%   fitd:action [ fitd:actionName "Skirmish" ; fitd:dots 4 ] ;
%   fitd:action [ fitd:actionName "Prowl" ; fitd:dots 3 ] .
% ════════════════════════════════════════════════════════════

\pgheader[SpinolaRed]{Ser Drago detto il Silenzioso}{Il Braccio / Il Problema Risolto}{Spinola}
\pgquote{Il mio nome non è importante. Il risultato sì.}

\section*{Background \& Motivazione}

\begin{rulebox}{}
  \textbf{Background:} Quarant'anni circa — non è sicuro nemmeno lui.
  Dalmata di nascita, mercenario per vocazione, uomo degli Spinola da
  otto anni. Non è un assassino — è un risolutore di problemi che a volte
  i problemi richiedono di risolvere con la violenza. Raffaele lo usa per
  le situazioni in cui i diplomatici non bastano: recupero crediti pericolosi,
  scorte sensibili, neutralizzazione di minacce fisiche. È l'unico membro
  della famiglia con vera competenza militare, e tutti se ne dimenticano
  finché non serve.

  \medskip\noindent\textcolor{SpinolaRed}{\textbf{Motivazione:}} Drago
  non ha posizioni politiche. Lavora per chi lo paga. Lavora per gli
  Spinola da abbastanza tempo da aver accumulato abbastanza per ritirarsi
  — ma non si è ancora ritirato. Non sa esattamente perché.
\end{rulebox}

\section*{Attributi \& Azioni}

\begin{multicols}{2}

\begin{attributoblock}[SpinolaRed]{INSIGHT --- Mente}
  \azione{Caccia}{Hunt}{3}{Trova persone, valuta minacce, pianifica azioni tattiche}
  \azione{Studio}{Study}{1}{Lingue (dalmatico, italiano, un po' di turco)}
  \azione{Osserva}{Survey}{2}{Legge le situazioni in termini di rischio fisico}
  \azione{Ingegno}{Tinker}{2}{Armi, trappole, ambienti modificati per la violenza}
\end{attributoblock}

\columnbreak

\begin{attributoblock}[SpinolaRed]{PROWESS --- Corpo}
  \azione{Finezza}{Finesse}{2}{Lama corta, lavoro preciso}
  \azione{Ombra}{Prowl}{3}{Movimento silenzioso, posizionamento tattico}
  \azione{Combattimento}{Skirmish}{4}{Professionale. Efficiente. Senza ornamenti.}
  \azione{Forza bruta}{Wreck}{3}{Sfonda, disabilita, neutralizza}
\end{attributoblock}

\end{multicols}

\begin{attributoblock}[SpinolaRed]{RESOLVE --- Spirito}
  \begin{multicols}{2}
    \azione{Percezione}{Attune}{2}{Sente quando una situazione sta per degenerare}
    \azione{Comando}{Command}{1}{Non comanda — esegue e protegge}
    \azione{Relazioni}{Consort}{1}{Mercenari, bassifondi, persone che non fanno domande}
    \azione{Persuasione}{Sway}{0}{Non è il suo metodo}
  \end{multicols}
\end{attributoblock}

\medskip
\begin{pgclockbox}{Stress \hfill \clock{9}{0}}
  \small\textit{Traumi: $\square$ Freddo \quad $\square$ Duro \quad $\square$ Ossessionato \quad $\square$ Paranoico \quad $\square$ Irrazionale \quad $\square$ Spezzato}
\end{pgclockbox}

\section*{Abilità Speciale}

% @rpg-schema: <.../pg-ser-drago>
%   fitd:specialAbility [ rdfs:label "Problema Risolto" ] .

\begin{spinolabox}{Problema Risolto}
  Quando Drago esegue un'azione violenta o coercitiva, può scegliere di farlo
  senza lasciare tracce fisiche o testimoniali riconducibili agli Spinola.
  Il problema sparisce. Le conseguenze narrative esistono ma non toccano
  la reputazione della famiglia.

  \medskip\noindent\textit{Una volta per colpo. Se Drago subisce Conseguenze
  Gravi nell'esecuzione, questa protezione decade — lui esiste, e parla.}
\end{spinolabox}

\section*{Legami}

\begin{table}[h]
\centering\renewcommand{\arraystretch}{1.4}
\begin{tabularx}{\linewidth}{@{}L{2.6cm} X@{}}
  \toprule\rowcolor{SpinolaRed}
  \textcolor{white}{\textbf{PG}} & \textcolor{white}{\textbf{Legame}} \\
  \midrule
  Raffaele & Mi paga. Non mi chiede mai come ho risolto i problemi.
            Apprezzo questa qualità in un datore di lavoro. \\
  \rowcolor{SpinolaRed!8}
  Lucrezia & L'ho protetta una volta senza che me lo chiedesse.
            Da allora mi guarda come se dovesse qualcosa. Non deve niente. \\
  Margherita & Strana alleata. Sa muoversi in posti dove io sono troppo
              visibile. Ci siamo capiti senza parole. \\
  \rowcolor{SpinolaRed!8}
  Ippolito & Mi preoccupa. Ha il coraggio di chi non ha ancora visto
            abbastanza da capire il costo reale delle cose. \\
  \bottomrule
\end{tabularx}
\end{table}

\section*{Equipaggiamento}
\begin{goldlist}
  \item Spada corta e coltello — nessuna marca, nessun ornamento
  \item Cotta di maglia leggera sotto gli abiti borghesi
  \item Kit per cambiare aspetto rapidamente: tre set di abiti diversi, barba posticcia
  \item Denaro sufficiente per corrompere una guardia o pagare un passaggio su una barca
\end{goldlist}

\clearpage

\begin{secretbox}
  \textbf{Segreto:} Drago ha ricevuto un'offerta da una terza parte — non
  francesi, non genovesi — per consegnare Raffaele Spinola vivo entro la
  fine della notte. L'offerta è economicamente irresistibile. Non l'ha ancora
  accettata. Non sa se è perché ha una coscienza o perché non si fida
  dell'offerente. La differenza, per lui, è rilevante.

  \medskip\noindent\textcolor{SecretRed}{\textbf{Agenda:}}
  \textit{Capire chi c'è dietro l'offerta prima di decidere cosa farne.
  Se la fonte è affidabile, valuterà. Se non lo è, consegnerà l'informazione
  a Raffaele — e farà pagare bene quella lealtà improvvisata.}
\end{secretbox}
